% Part: first-order-logic
% Chapter: axiomatic-deduction
% Section: rules-and-proofs

\documentclass[../../../include/open-logic-section]{subfiles}

\begin{document}

\iftag{FOL}
      {\olfileid{fol}{axd}{rul}}
      {\olfileid{pl}{axd}{rul}}

\olsection{नियम आणि \usetoken{P}{derivation}}

\begin{explain}
  स्वयंसिद्धकीय !!{derivation}s ही तर्कशास्त्रासाठी कदाचित सर्वांत सोपी
  !!{derivation} पद्धत आहे. !!^a{derivation} म्हणजे केवळ !!{formula}s चा
  अनुक्रम. !!a{derivation} म्हणून गणले जाण्यासाठी अनुक्रमातील प्रत्येक
  !!{formula} एकतर स्वयंसिद्धकाचे उदाहरण असले पाहिजे, किंवा अनुक्रमात
  त्याच्या आधी येणाऱ्या एक वा अधिक !!{formula}s पासून निगमन नियमाने
  निष्पन्न झाले पाहिजे. !!^a{derivation} तिचे शेवटचे !!{formula}
  !!{derive}s करते.
\end{explain}

\begin{defn}[!!^{derivability}]
जर $\Gamma$ हा $\Lang L$ मधील !!{formula}s चा संच असेल, तर
$\Gamma$ पासूनची \emph{!!{derivation}} म्हणजे !!{formula}s चा सांत अनुक्रम
$!A_1$, \dots,~$!A_n$ असा की प्रत्येक $i \le n$ साठी पुढीलपैकी एक अट
पूर्ण होते:
\begin{enumerate}
\item $!A_i \in \Gamma$; किंवा
\item $!A_i$ हे स्वयंसिद्धक आहे; किंवा
\item $j < i$ (आणि $k < i$) असलेल्या एखाद्या $!A_j$ (आणि $!A_k$)
  पासून निगमन नियमाने $!A_i$ निष्पन्न होते.
\end{enumerate}
\end{defn}

योग्य !!{derivation} म्हणून कशाची गणना होते हे आपण कोणते निगमन नियम
मान्य करतो (आणि अर्थातच कोणती सूत्रे स्वयंसिद्धके मानतो) यावर अवलंबून
असते. निगमन नियम म्हणजे असे जर-तर विधान जे विशिष्ट अटींमध्ये
!!a{derivation} मधील पायरी~$A_i$ ही निगमनाची योग्य पायरी आहे असे सांगते.

\begin{defn}[निगमन नियम]
\emph{निगमन नियम} हा $\Gamma$ पासूनच्या !!a{derivation} मध्ये निगमनाची
योग्य पायरी म्हणून कशाची गणना होते यासाठी पुरेशी अट देतो.
\end{defn}

उदाहरणार्थ, $!A \in \Gamma$ असलेला कोणताही एक-घटकी अनुक्रम $!A$ हा
आपोआपच !!a{derivation} म्हणून गणला जात असल्यामुळे, पुढील नियम हा
अतिशय सोपा निगमन नियम असू शकतो:
\begin{quote}
  जर $!A \in \Gamma$, तर $\Gamma$ पासूनच्या कोणत्याही !!{derivation}
  मध्ये $!A$ ही नेहमीच निगमनाची योग्य पायरी असते.
\end{quote}
त्याचप्रमाणे, $!A$ हे स्वयंसिद्धकांपैकी एक असेल, तर एकटे $!A$ ही
!!a{derivation} असते; म्हणून पुढील विधानही निगमन नियम आहे:
\begin{quote}
  जर $!A$ हे स्वयंसिद्धक असेल, तर $!A$ ही निगमनाची योग्य पायरी असते.
\end{quote}
विचाराधीन पायरीच्या आधी दिसणाऱ्या !!{formula}s चा निगमन नियम आधार घेतो
तेव्हा गोष्ट अधिक रोचक होते. पुढील नियमाला \emph{विध्यात्मक अनुमान}
म्हणतात:
\begin{quote}
  जर $!B \lif !A$ आणि $!B$ हे !!{derivation} मध्ये यापूर्वी आढळत असतील,
  तर~$!A$ ही निगमनाची योग्य पायरी असते.
\end{quote}
हा एकमेव निगमन नियम असेल, तर वर दिलेल्या !!{derivation} च्या व्याख्येचा
अर्थ असा होतो: प्रत्येक $i \le n$ साठी पुढीलपैकी एक अट पूर्ण होत असेल
तेव्हा आणि केवळ तेव्हाच $!A_1$, \dots,~$!A_n$ ही !!a{derivation} असते:
\begin{enumerate}
\item $!A_i \in \Gamma$; किंवा
\item $!A_i$ हे स्वयंसिद्धक आहे; किंवा
\item एखाद्या $j < i$ साठी $!A_j$ हे $!B \lif !A_i$ आहे, आणि एखाद्या
  $k < i$ साठी $!A_k$ हे~$!B$ आहे.
\end{enumerate}
शेवटचे कलम असे सांगते की~$!A_j$ ($!B \lif !A_i$) आणि $!A_k$ ($!B$)
यांपासून विध्यात्मक अनुमानाने $!A_i$ निष्पन्न होते. आपण $1$ पासून~$n$
पर्यंत जाऊ शकलो आणि प्रत्येक वेळी~$\Gamma$ मध्ये असलेले, स्वयंसिद्धक
असलेले किंवा निगमन नियमाने योग्य निगमन-पायरी ठरणारे !!a{formula}~$!A_i$
मिळाले, तर संपूर्ण अनुक्रम योग्य !!{derivation} म्हणून गणला जातो.

\begin{defn}[!!^{derivability}]
जर $\Gamma$ पासून सुरू होऊन $!A$ वर संपणारी !!a{derivation} असेल, तर
!!{formula}~$!A$ हे $\Gamma$ पासून \emph{!!{derivable}} असते; हे आपण
$\Gamma \Proves !A$ असे लिहितो.
\end{defn}

\begin{defn}[प्रमेये]
रिक्त संचापासून $!A$ ची !!a{derivation} असेल, तर !!^a{formula}~$!A$ हे
\emph{प्रमेय} असते. $!A$ हे प्रमेय असेल तर आपण $\Proves !A$ लिहितो आणि
ते प्रमेय नसेल तर $\Proves/ !A$ लिहितो.
\end{defn}


\end{document}
